% Author response template for the reporting checklist
% Standalone usage: pdflatex author_response_template.tex
\documentclass[10pt,a4paper]{article}

\usepackage[margin=20mm]{geometry}
\usepackage[T1]{fontenc}
\usepackage[utf8]{inputenc}
\usepackage{array}
\usepackage{booktabs}
\usepackage{longtable}
\usepackage{ragged2e}
\usepackage[hidelinks]{hyperref}

\newcolumntype{L}[1]{>{\RaggedRight\arraybackslash}p{#1}}
\newcolumntype{N}[1]{>{\RaggedRight\hyphenpenalty=10000\exhyphenpenalty=10000\arraybackslash}p{#1}}
\setlength{\tabcolsep}{3pt}
\renewcommand{\arraystretch}{1.12}

\newcommand{\ResponseTableHeader}{%
\toprule
\textbf{ID} & \textbf{Item} & \textbf{State} & \textbf{Response} & \textbf{Source location / artefact}\\
\midrule}

\begin{document}

\section*{Author response to the reporting checklist}

\noindent\textbf{Paper title:} \emph{(fill in)}\\
\textbf{Version and date:} \emph{(fill in)}\\
\textbf{Corresponding author:} \emph{(fill in)}\\
\textbf{Code or data release:} \emph{(URL, DOI, or repository tag/commit)}

\medskip
\noindent\textbf{Declared workflow, cohort and endpoint:} \emph{(candidate population, physical evidence level and comparison)}\\
\textbf{Development format:} \texttt{1.0.0-draft}. The archived 0.2.0 records retain their original interpretation.

\medskip
\noindent\textbf{Instructions.} Supply operational values and decisions with exact source locations, rather than a list of states alone. Distinguish \texttt{reported}, mixed \texttt{partially\_reported}, unavailable \texttt{not\_reported}, known \texttt{not\_performed}, and justified \texttt{not\_applicable}. Explain the reason and claim consequence of non-performance and the rationale for inapplicability. Keep unavailable quantities missing rather than zero. Attach the consolidated record to the submission and refer to it briefly in the main text. The identifiers below are retained for machine compatibility, not a requirement to print numerical category codes in the manuscript.

\setlength{\LTpre}{4pt}
\setlength{\LTpost}{8pt}

\subsection*{Core response categories}

\begin{longtable}{@{}L{0.05\textwidth} N{0.16\textwidth} L{0.14\textwidth} L{0.37\textwidth} L{0.20\textwidth}@{}}
\ResponseTableHeader
\endfirsthead
\ResponseTableHeader
\endhead
\bottomrule
\endfoot
\texttt{D1} & Data & \emph{(state)} & \emph{(response)} & \emph{(exact location)}\\
\texttt{S1} & Splits & \emph{(state)} & \emph{(response)} & \emph{(exact location)}\\
\texttt{M1} & Inputs/Model & \emph{(state)} & \emph{(response)} & \emph{(exact location)}\\
\texttt{T1} & Training & \emph{(state)} & \emph{(response)} & \emph{(exact location)}\\
\texttt{G1} & Generation & \emph{(state)} & \emph{(response)} & \emph{(exact location)}\\
\texttt{E1} & Metrics & \emph{(state)} & \emph{(response)} & \emph{(exact location)}\\
\texttt{B1} & Baselines and ablations & \emph{(state)} & \emph{(response)} & \emph{(exact location)}\\
\texttt{F1} & Feasibility & \emph{(state)} & \emph{(response)} & \emph{(exact location)}\\
\texttt{R1} & Reproducibility & \emph{(state)} & \emph{(response)} & \emph{(exact location)}\\
\end{longtable}

\subsection*{Context-dependent reporting items}

\begin{longtable}{@{}L{0.05\textwidth} N{0.16\textwidth} L{0.14\textwidth} L{0.37\textwidth} L{0.20\textwidth}@{}}
\ResponseTableHeader
\endfirsthead
\ResponseTableHeader
\endhead
\bottomrule
\endfoot
\texttt{CL1} & Closed-loop/UQ & \emph{(state)} & \emph{(response)} & \emph{(exact location)}\\
\texttt{C1} & Compute footprint & \emph{(state)} & \emph{(response)} & \emph{(exact location)}\\
\texttt{GOV1} & Governance and ethics & \emph{(state)} & \emph{(response)} & \emph{(exact location)}\\
\end{longtable}

\subsection*{Optional key-number summary}

\begin{longtable}{@{}L{0.05\textwidth} N{0.16\textwidth} L{0.14\textwidth} L{0.37\textwidth} L{0.20\textwidth}@{}}
\ResponseTableHeader
\endfirsthead
\ResponseTableHeader
\endhead
\bottomrule
\endfoot
\texttt{K1} & Key-number summary & \emph{(state)} & \emph{(response)} & \emph{(exact location)}\\
\end{longtable}

\noindent\footnotesize\textbf{Suggested K1 fields.} Total candidates; stagewise pass rates; validated outcomes; evaluator budget; closed-loop rounds and evaluations, where used; training and inference cost; and the versioned code or record identifier.

\medskip\noindent The compact summary is optional as a format; claim-relevant counts remain required. Separate candidate counts from evaluator calls, failures and repeats in the stage ledger. Report inference throughput under Generation and training time under Training.

\end{document}
